\section{Results}
\label{sec:results}
\subsection{Accuracy and Calibration at a Glance}
Table~\ref{tab:main} compares DS2S-E with S2S-E, the forecaster of our earlier deployment, on the RTT of the test sessions. The interval score improves on all six links, by \BSISWorstAbs\% to \BSISBestAbs\%, and the session-bootstrap interval of every change excludes zero. DS2S-E has the lower interval score in each of the \BSnSessTot{} test sessions. Both forecasters keep their marginal coverage close to 95\%, lowest on Starlink UL for both. Their difference therefore lies in the width of the intervals, not in how often they cover. Fig.~\ref{fig:heat} shows that the gain holds at every horizon step. It is largest at short horizons on the satellite links and flat on 5G.

\begin{table*}[t]
\centering
\caption{RTT on the test sessions, averaged over 0.2--3\,s, for the three-seed ensembles S2S-E and DS2S-E. Changes of DS2S-E against S2S-E with 95\% session-bootstrap intervals (2000 replicates). Bold: the interval excludes zero. MAE uses the median of DS2S-E, RMSE the means. IS and PICP refer to the calibrated 95\% interval.}
\label{tab:main}
\footnotesize
\setlength{\tabcolsep}{2.9pt}
\begin{tabular}{l rrr rrr rrr cc}
\toprule
 & \multicolumn{3}{c}{Interval score (ms)} & \multicolumn{3}{c}{MAE (ms)} & \multicolumn{3}{c}{RMSE (ms)} & \multicolumn{2}{c}{PICP} \\
\cmidrule(lr){2-4}\cmidrule(lr){5-7}\cmidrule(lr){8-10}\cmidrule(lr){11-12}
Link & S2S-E & DS2S-E & $\Delta$\% [95\% CI] & S2S-E & DS2S-E & $\Delta$\% [95\% CI] & S2S-E & DS2S-E & $\Delta$\% [95\% CI] & S2S-E & DS2S-E \\
\midrule
\tabinput{figs/data/tab_main_boot.tex}
\bottomrule
\end{tabular}
\end{table*}

\begin{figure}[t]
  \centering
  % fr_heat_col.tex - generated by modeling/experiments/fr_figures.py, do not edit by hand.
\begin{tikzpicture}
\begin{groupplot}[group style={group size=1 by 2, vertical sep=0.95cm, xticklabels at=edge bottom, xlabels at=edge bottom},
  width=6.0cm, height=3.55cm, enlargelimits=false, xmin=0.5, xmax=15.5, ymin=-0.5, ymax=5.5, ytick={0,...,5}, yticklabels={Starlink DL,OneWeb DL,5G UL,Starlink UL,OneWeb UL,5G DL},
  xtick={1,5,10,15}, xticklabels={0.2,1,2,3}, xlabel={horizon (s)}, tick label style={font=\scriptsize}, label style={font=\scriptsize},
  title style={font=\footnotesize, yshift=-0.6ex}, colormap={bwr}{rgb255=(33,102,172) rgb255=(247,247,247) rgb255=(178,24,43)}, point meta min=-50.0, point meta max=50.0, y dir=reverse, clip=false]
\nextgroupplot[title={(a) RTT MAE change (\%)}]
\addplot[matrix plot*, mesh/cols=15, point meta=explicit] table[x=x, y=y, meta=v] {figs/data/fr_heat_MAE.dat};
\node[anchor=west, font=\scriptsize] at (axis cs:15.6,0) {$-15$};
\node[anchor=west, font=\scriptsize] at (axis cs:15.6,1) {$-10$};
\node[anchor=west, font=\scriptsize] at (axis cs:15.6,2) {$-18$};
\node[anchor=west, font=\scriptsize] at (axis cs:15.6,3) {$-3$};
\node[anchor=west, font=\scriptsize] at (axis cs:15.6,4) {$-2$};
\node[anchor=west, font=\scriptsize] at (axis cs:15.6,5) {$-2$};
\nextgroupplot[title={(b) RTT interval-score change (\%)}, colorbar horizontal, colorbar style={at={(0.5,-0.62)}, anchor=north, height=0.16cm, xtick={-50,-25,0,25,50}, xticklabel style={font=\scriptsize}}]
\addplot[matrix plot*, mesh/cols=15, point meta=explicit] table[x=x, y=y, meta=v] {figs/data/fr_heat_MIS.dat};
\node[anchor=west, font=\scriptsize] at (axis cs:15.6,0) {$-41$};
\node[anchor=west, font=\scriptsize] at (axis cs:15.6,1) {$-20$};
\node[anchor=west, font=\scriptsize] at (axis cs:15.6,2) {$-28$};
\node[anchor=west, font=\scriptsize] at (axis cs:15.6,3) {$-9$};
\node[anchor=west, font=\scriptsize] at (axis cs:15.6,4) {$-26$};
\node[anchor=west, font=\scriptsize] at (axis cs:15.6,5) {$-14$};
\end{groupplot}
\end{tikzpicture}

  \caption{Change (\%) of the RTT MAE and of the calibrated 95\% interval score of DS2S-E against S2S-E per link and horizon step. Blue: DS2S-E is better. Right of each row: the change over all steps.}
  \label{fig:heat}
\end{figure}

The point accuracy changes much less. The MAE improves significantly on five links, and the RMSE changes by at most \BSRMSEMaxAbs\%. This is what Section~\ref{sec:predictable} leads one to expect, since the level of the next RTT is hard to forecast. Almost all of the MAE gain comes from reporting the median instead of the mean. On the five links other than Starlink DL, the mean of DS2S-E changes the MAE by \MeanMAELo\% to $+$\MeanMAEHi\%, against \MedMAELo\% to \MedMAEHi\% for its median. Under right-skewed errors the median is closer to a typical outcome, and by Proposition~\ref{prop:optimal} it is the optimal forecast for the absolute error. Starlink DL is the one link where the mean also improves (\MeanMAESLDL\%), because its sawtooth is partly predictable.

\subsection{Why the Interval Improves}
Fig.~\ref{fig:condcov} breaks the coverage down by the volatility of the input window and the horizon. The constant interval of S2S-E is blue in the calm deciles and red in the volatile ones. It over-covers quiet windows and under-covers volatile ones. It covers the calmest tenth of the windows at \CovConstSLDLdOne{} and the most volatile tenth at \CovConstSLDLdTen{} on Starlink DL, and at \CovConstQULdOne{} and \CovConstQULdTen{} on 5G UL. Scaling the interval in proportion to the input SD, with a small floor, reverses the error (bottom row). The future error grows much more slowly than the input SD, and the calmest windows are now under-covered (OneWeb DL: \CovNormOWDLdOne). The DS2S-E interval is close to white on Starlink DL and 5G UL. Its coverage gap is \GapDSSLDL{} on Starlink DL, \GapDSQDL{} on 5G DL, and \GapDSQUL{} on 5G UL, against \GapConstSLDL{}, \GapConstQDL{}, and \GapConstQUL{} for the constant interval. The reduction is significant on these three links and on Starlink UL. On Starlink DL the constant interval is \WidConstSLDL\,ms wide at 1\,s ahead in every window. The DS2S-E interval is \WidDSSLDLcalmLo{} to \WidDSSLDLcalmHi\,ms wide in the three calmest deciles, where the link is idle, and \WidDSSLDLloadLo{} to \WidDSSLDLloadHi\,ms in the loaded ones.

\begin{figure*}[t]
  \centering
  % fr_condcov.tex - generated by modeling/experiments/fr_condcov.py, do not edit by hand.
% Conditional coverage of the calibrated 95% RTT bands (figure*, full text width).
% Cell = share of TEST windows whose RTT at that horizon step fell inside the band, among the
% windows whose 6 s input lies in that input-volatility decile (edges from validation windows).
% White = nominal 0.95, blue = over-coverage (band wider than needed), red = under-coverage.
% Values below 0.85 share the darkest red.
\begin{tikzpicture}
\begin{groupplot}[
  group style={group size=3 by 3, horizontal sep=0.28cm, vertical sep=0.62cm,
    xticklabels at=edge bottom, yticklabels at=edge left, xlabels at=edge bottom, ylabels at=edge left},
  scale only axis, width=4.75cm, height=2.2cm, enlargelimits=false, axis on top,
  xmin=0.5, xmax=15.5, ymin=0.5, ymax=10.5,
  xtick={1,5,10,15}, xticklabels={0.2,1,2,3}, ytick={1,5,10}, minor ytick={2,3,4,6,7,8,9},
  tick label style={font=\scriptsize}, label style={font=\scriptsize},
  title style={font=\footnotesize, yshift=-0.9ex}, ylabel style={align=center, yshift=-0.5ex},
  xlabel={horizon (s)},
  colormap={covdiv}{rgb255(0cm)=(103,0,31); rgb255(1cm)=(178,24,43); rgb255(1.5cm)=(239,138,98);
    rgb255(2cm)=(247,247,247); rgb255(2.5cm)=(146,197,222); rgb255(3cm)=(33,102,172)},
  point meta min=0.85, point meta max=1.0,
]
\nextgroupplot[title={(a) Starlink DL}, ylabel={input-SD decile}, colorbar horizontal, colorbar to name=cbcondcov, colorbar style={width=6.2cm, height=0.2cm, xtick={0.85,0.9,0.95,1}, tick label style={font=\scriptsize}, xticklabel style={/pgf/number format/fixed, /pgf/number format/fixed zerofill, /pgf/number format/precision=2}}]
\addplot[matrix plot*, mesh/cols=15, point meta=explicit] table[x=h, y=dec, meta=covs2sconst] {figs/data/fr_condcov_heat_starlink_dl.dat};
\nextgroupplot[title={(b) OneWeb DL}]
\addplot[matrix plot*, mesh/cols=15, point meta=explicit] table[x=h, y=dec, meta=covs2sconst] {figs/data/fr_condcov_heat_oneweb_dl.dat};
\nextgroupplot[title={(c) 5G UL}]
\addplot[matrix plot*, mesh/cols=15, point meta=explicit] table[x=h, y=dec, meta=covs2sconst] {figs/data/fr_condcov_heat_quectel_ul.dat};
\nextgroupplot[title={(d) Starlink DL}, ylabel={input-SD decile}]
\addplot[matrix plot*, mesh/cols=15, point meta=explicit] table[x=h, y=dec, meta=covds2scqr] {figs/data/fr_condcov_heat_starlink_dl.dat};
\nextgroupplot[title={(e) OneWeb DL}]
\addplot[matrix plot*, mesh/cols=15, point meta=explicit] table[x=h, y=dec, meta=covds2scqr] {figs/data/fr_condcov_heat_oneweb_dl.dat};
\nextgroupplot[title={(f) 5G UL}]
\addplot[matrix plot*, mesh/cols=15, point meta=explicit] table[x=h, y=dec, meta=covds2scqr] {figs/data/fr_condcov_heat_quectel_ul.dat};
\nextgroupplot[title={(g) Starlink DL}, ylabel={input-SD decile}]
\addplot[matrix plot*, mesh/cols=15, point meta=explicit] table[x=h, y=dec, meta=covs2snorm] {figs/data/fr_condcov_heat_starlink_dl.dat};
\nextgroupplot[title={(h) OneWeb DL}]
\addplot[matrix plot*, mesh/cols=15, point meta=explicit] table[x=h, y=dec, meta=covs2snorm] {figs/data/fr_condcov_heat_oneweb_dl.dat};
\nextgroupplot[title={(i) 5G UL}]
\addplot[matrix plot*, mesh/cols=15, point meta=explicit] table[x=h, y=dec, meta=covs2snorm] {figs/data/fr_condcov_heat_quectel_ul.dat};
\end{groupplot}
\node[rotate=90, anchor=south, align=center, text width=2.4cm, font=\scriptsize] at ($(group c1r1.west)+(-1.05cm,0)$) {\textbf{S2S-E}\\constant band};
\node[rotate=90, anchor=south, align=center, text width=2.4cm, font=\scriptsize] at ($(group c1r2.west)+(-1.05cm,0)$) {\textbf{DS2S-E}\\CQR band};
\node[rotate=90, anchor=south, align=center, text width=2.4cm, font=\scriptsize] at ($(group c1r3.west)+(-1.05cm,0)$) {\textbf{S2S-E}\\SD-scaled band};
\node[draw=black!50, fill=white, inner sep=3pt, anchor=south] at ($(group c2r1.north)+(0,0.62cm)$)
  {\scriptsize coverage of the nominal 95\% band\quad\raisebox{-0.5ex}{\pgfplotscolorbarfromname{cbcondcov}}\quad
   \begin{tabular}{@{}l@{}}red: band too narrow\\[-1pt] blue: band too wide\end{tabular}};
\end{tikzpicture}

  \caption{Conditional coverage of the calibrated 95\% RTT intervals on the test sessions. A cell is the coverage among the windows whose input SD lies in that decile (edges from the calibration sessions) at that horizon step. White: nominal 0.95. Blue: the interval is wider than needed. Red: it is too narrow. Values below 0.85 share the darkest red. Bottom row: interval scaled by the input SD with a small floor.}
  \label{fig:condcov}
\end{figure*}

The OneWeb links show the limit of DS2S-E. Its coverage gap there is not smaller than that of the constant interval. DS2S-E widens its interval only \WidRatioOWDL{} times (DL) and \WidRatioOWUL{} times (UL) from the calmest to the most volatile tenth, while the calibrated error grows \ErrRatioOWDL{} and \ErrRatioOWUL{} times, and it covers that tenth only at \CovDSOWDLdTen{} and \CovDSOWULdTen{}. The input carries the warning, since an SD-scaled interval covers the same tenth at \CovNormTunedOWDLdTen{} and \CovNormTunedOWULdTen. DS2S-E under-reacts to the rare OneWeb bursts.

Where does the gain in interval score come from? We built intermediate intervals that change one ingredient at a time, all calibrated on the same sessions. Moving a constant interval from the S2S-E mean to the DS2S-E mean changes the score by \DecCentreLo\% to $+$\DecCentreHi\%. The better center is not the source. Giving the S2S-E forecast the per-window width of DS2S-E, placed symmetrically, gains \DecWidthLoAbs\% to \DecWidthHiAbs\%. Placing the interval asymmetrically, as DS2S-E does, gains another \DecAsymLoAbs\% to \DecAsymHiAbs\%. The adaptive width carries most of the gain on Starlink DL and OneWeb UL. The skew carries most of it on 5G UL, where the raw DS2S-E interval at 1\,s reaches \AsymQULratio{} times further above its median than below it.

\subsection{How Specific Is the Gain to DS2S?}
Table~\ref{tab:extent} lists the interval score and the coverage gap of every interval we built. It bounds how much of the gain requires our model. First, no interval around the S2S-E point forecast comes close to the quantile forecasters. The asymmetric interval helps where the skew is strong and the SD-scaled one where the regime shows in the input SD, but the best of them is still \ExtBestPointLoAbs\% to \ExtBestPointHiAbs\% worse than DS2S-E on every link. The SD-scaled interval balances the decile coverage well on OneWeb because it is scaled by the variable that defines the deciles, and it pays for this with a higher interval score. Second, the new inputs and the anchored decoder do not explain the gain. S2S-AE-E shares them with DS2S-E but outputs a mean, and it changes the interval score of S2S-E by only \ExtFrfullLo\% to $+$\ExtFrfullHi\%. \GaussSentence{} Fourth, the linear quantile model LQR does almost as well as DS2S-E. It is within \ExtLQRMaxAbs\% of DS2S-E on every link and better on OneWeb UL and Starlink UL, and DS2S-E is significantly better only on Starlink DL, 5G DL, and 5G UL, by \LQRsigLo\% to \LQRsigHi\%.

\begin{table*}[t]
\centering
\caption{Interval score (IS, ms) and coverage gap $\max_d|\mathrm{cov}_d-0.95|$ over input-SD deciles of all 95\% RTT intervals on the test sessions. All intervals are calibrated on the calibration sessions. SD-scaled: floor tuned on the calibration sessions. S2S-AE-E: the inputs and decoder of DS2S-E with a scalar output. Bold: lowest IS of the link. GS2S-E \GaussSeedNote.}
\label{tab:extent}
\footnotesize
\setlength{\tabcolsep}{2.7pt}
\begin{tabular}{ll*{12}{r}}
\toprule
 & & \multicolumn{2}{c}{OneWeb DL} & \multicolumn{2}{c}{OneWeb UL} & \multicolumn{2}{c}{Starlink DL} & \multicolumn{2}{c}{Starlink UL} & \multicolumn{2}{c}{5G DL} & \multicolumn{2}{c}{5G UL} \\
\cmidrule(lr){3-4}\cmidrule(lr){5-6}\cmidrule(lr){7-8}\cmidrule(lr){9-10}\cmidrule(lr){11-12}\cmidrule(lr){13-14}
Forecast & Interval & IS & gap & IS & gap & IS & gap & IS & gap & IS & gap & IS & gap \\
\midrule
\tabinput{figs/data/tab_extent.tex}
\bottomrule
\end{tabular}
\end{table*}

The source of the gain is therefore the distributional output calibrated across sessions, which a linear model can exploit too. What the recurrent network adds is a better point forecast. Against LQR, DS2S-E lowers the RMSE on all six links (\LQRRMSELoAbs\% to \LQRRMSEHiAbs\%, all significant), and it lowers the interval score further on the three links where the RTT has the most structure.

\subsection{Can the Results Be Trusted?}
Four checks support these results. First, the session is the unit of evidence. A bootstrap over windows treats the \BSnWinTot{} overlapping test windows as independent. It gives intervals \BSWidRatioLo{} to \BSWidRatioHi{} times narrower than the session bootstrap and would call two inconclusive changes significant, the RMSE on OneWeb DL and the MAE on OneWeb UL. Second, no single session drives a result. Leaving out any one test session moves the interval-score change by at most \BSLooMax{} points. Third, the gain is visible only with calibration on held-out sessions. With intervals from the empirical quantiles of the test residuals, a common practice that yields 95\% coverage by construction, DS2S-E and S2S-E are within \IsProtoMaxAbs\% of each other on every link, because both intervals then have a constant width. Fourth, the inputs are causal (Section~\ref{sec:causal}), every tuning choice uses the calibration sessions only, and on Starlink DL the comparison holds on the measured targets alone (MAE \RealMAEFRSLDL{} against \RealMAEPrevSLDL\,ms).

\subsection{Value for a Deadline-Aware Scheduler}
We emulate deadline-aware packet duplication on the test sessions of the four satellite links, with the satellite path as primary and the 5G path of the same direction as backup. At time $t$ the scheduler plans the packets that it will send at $t+h$, $h\in\{0.2,\dots,3\}$\,s. A packet on the primary misses its deadline $D$ if the measured RTT at $t+h$ exceeds $D$, which we set to the 90th percentile of the calibration-session RTT (\SchedDSLDL\,ms on Starlink DL). The scheduler duplicates a packet on the backup when its forecast bound exceeds $D$. We count the share of duplicated packets, which costs backup capacity, and the share of packets that are late and were not duplicated (misses). The backup meets the downlink deadlines in all but \BackupDLmax\% of the 5G DL samples. The 5G UL exceeds the Starlink UL deadline in \BackupQULSL\% of its samples. If the two paths are independent, this raises the miss rate of the S2S-E bound on Starlink UL from \BackupSSmissSLUL\% to \BackupSSmissCorrSLUL\% and that of DS2S-E from \BackupDSmissSLUL\% to \BackupDSmissCorrSLUL\%.

\begin{figure*}[t]
  \centering
  \resizebox{\textwidth}{!}{% fr_scheduler.tex - generated by modeling/experiments/fr_scheduler.py (edit the script, not this file).
% Deadline-aware duplication on the test sessions: duplication rate vs deadline-miss rate, D = validation p90.
\pgfplotsset{
  schOracle/.style={black, dashed, thick, mark=none},
  schRandom/.style={black!45, thin, mark=none},
  schReact/.style={cPers, densely dotted, thick, mark=x, mark size=1.7pt, mark options={solid}, mark repeat=20, mark phase=11},
  schS2S/.style={mPrev, mark repeat=20, mark phase=11},
  schAdapt/.style={cResid, thick, dashed, mark=triangle, mark size=1.6pt, mark options={solid}, mark repeat=20, mark phase=11},
  schDS2S/.style={mFR, mark repeat=20, mark phase=11},
  schGauss/.style={cRidge, thick, dashdotted, mark=diamond, mark size=1.6pt, mark options={solid}, mark repeat=20, mark phase=11},
  schCalS2S/.style={only marks, mark=triangle*, mark size=2.4pt, mark options={draw=black, fill=cPrev, thin}},
  schCalDS2S/.style={only marks, mark=pentagon*, mark size=2.4pt, mark options={draw=black, fill=cFR, thin}},
  schRing/.style={only marks, mark=o, mark size=3.9pt, black, semithick, forget plot},
  schInset/.style={scale only axis, width=1.6cm, height=1.15cm, axis background/.style={fill=white},
    tick label style={font=\tiny}, xtick pos=bottom, ytick pos=left, grid=major, grid style={gray!20},
    tick align=inside, major tick length=1.5pt, clip=true},
}
\begin{tikzpicture}
\begin{groupplot}[group style={group size=4 by 1, horizontal sep=1.15cm}, frbase, scale only axis,
  width=3.45cm, height=3.0cm, xmin=0, xmax=100, xtick={0,25,50,75,100}, ymin=0,
  xlabel={duplication rate (\%)}, title style={font=\footnotesize, yshift=-0.6ex}]
\nextgroupplot[title={(a) Starlink DL, $D=198$\,ms}, ymax=8, ylabel={deadline-miss rate (\%)}, legend to name=legsched, legend columns=5, legend cell align=left, legend style={font=\scriptsize, draw=black!60, /tikz/every even column/.append style={column sep=0.22cm}}]
\addplot[schOracle] table[x=dup, y=oracle] {figs/data/fr_scheduler_curve_starlink_dl.dat}; \addlegendentry{Oracle}
\addplot[schRandom] table[x=dup, y=random] {figs/data/fr_scheduler_curve_starlink_dl.dat}; \addlegendentry{Random}
\addplot[schReact] table[x=dup, y=reactive] {figs/data/fr_scheduler_curve_starlink_dl.dat}; \addlegendentry{Reactive}
\addplot[schS2S] table[x=dup, y=s2s] {figs/data/fr_scheduler_curve_starlink_dl.dat}; \addlegendentry{S2S-E}
\addplot[schAdapt] table[x=dup, y=s2sadapt] {figs/data/fr_scheduler_curve_starlink_dl.dat}; \addlegendentry{S2S-E adapt.}
\addplot[schGauss] table[x=dup, y=gauss] {figs/data/fr_scheduler_curve_starlink_dl.dat}; \addlegendentry{Gauss-E}
\addplot[schDS2S] table[x=dup, y=ds2s] {figs/data/fr_scheduler_curve_starlink_dl.dat}; \addlegendentry{DS2S-E}
\addplot[schCalS2S] table[x=s2sdup, y=s2smiss] {figs/data/fr_scheduler_cal_starlink_dl.dat}; \addlegendentry{S2S-E bound}
\addplot[schCalDS2S] table[x=ds2sdup, y=ds2smiss] {figs/data/fr_scheduler_cal_starlink_dl.dat}; \addlegendentry{DS2S-E bound}
\addplot[schRing, skip coords between index={0}{2}, skip coords between index={3}{5}] table[x=s2sdup, y=s2smiss] {figs/data/fr_scheduler_cal_starlink_dl.dat};
\addplot[schRing, skip coords between index={0}{2}, skip coords between index={3}{5}] table[x=ds2sdup, y=ds2smiss] {figs/data/fr_scheduler_cal_starlink_dl.dat};
\draw[black!65, thin] (axis cs:0,0.688) rectangle (axis cs:20,7.717);
\coordinate (zsrca) at (axis cs:20,4.203);
\nextgroupplot[title={(b) Starlink UL, $D=51$\,ms}, ymax=5]
\addplot[schOracle] table[x=dup, y=oracle] {figs/data/fr_scheduler_curve_starlink_ul.dat};
\addplot[schRandom] table[x=dup, y=random] {figs/data/fr_scheduler_curve_starlink_ul.dat};
\addplot[schReact] table[x=dup, y=reactive] {figs/data/fr_scheduler_curve_starlink_ul.dat};
\addplot[schS2S] table[x=dup, y=s2s] {figs/data/fr_scheduler_curve_starlink_ul.dat};
\addplot[schAdapt] table[x=dup, y=s2sadapt] {figs/data/fr_scheduler_curve_starlink_ul.dat};
\addplot[schGauss] table[x=dup, y=gauss] {figs/data/fr_scheduler_curve_starlink_ul.dat};
\addplot[schDS2S] table[x=dup, y=ds2s] {figs/data/fr_scheduler_curve_starlink_ul.dat};
\addplot[schCalS2S] table[x=s2sdup, y=s2smiss] {figs/data/fr_scheduler_cal_starlink_ul.dat};
\addplot[schCalDS2S] table[x=ds2sdup, y=ds2smiss] {figs/data/fr_scheduler_cal_starlink_ul.dat};
\addplot[schRing, skip coords between index={0}{2}, skip coords between index={3}{5}] table[x=s2sdup, y=s2smiss] {figs/data/fr_scheduler_cal_starlink_ul.dat};
\addplot[schRing, skip coords between index={0}{2}, skip coords between index={3}{5}] table[x=ds2sdup, y=ds2smiss] {figs/data/fr_scheduler_cal_starlink_ul.dat};
\draw[black!65, thin] (axis cs:2,1.663) rectangle (axis cs:20,4.212);
\coordinate (zsrcb) at (axis cs:20,2.937);
\nextgroupplot[title={(c) OneWeb DL, $D=98$\,ms}, ymax=6]
\addplot[schOracle] table[x=dup, y=oracle] {figs/data/fr_scheduler_curve_oneweb_dl.dat};
\addplot[schRandom] table[x=dup, y=random] {figs/data/fr_scheduler_curve_oneweb_dl.dat};
\addplot[schReact] table[x=dup, y=reactive] {figs/data/fr_scheduler_curve_oneweb_dl.dat};
\addplot[schS2S] table[x=dup, y=s2s] {figs/data/fr_scheduler_curve_oneweb_dl.dat};
\addplot[schAdapt] table[x=dup, y=s2sadapt] {figs/data/fr_scheduler_curve_oneweb_dl.dat};
\addplot[schGauss] table[x=dup, y=gauss] {figs/data/fr_scheduler_curve_oneweb_dl.dat};
\addplot[schDS2S] table[x=dup, y=ds2s] {figs/data/fr_scheduler_curve_oneweb_dl.dat};
\addplot[schCalS2S] table[x=s2sdup, y=s2smiss] {figs/data/fr_scheduler_cal_oneweb_dl.dat};
\addplot[schCalDS2S] table[x=ds2sdup, y=ds2smiss] {figs/data/fr_scheduler_cal_oneweb_dl.dat};
\addplot[schRing, skip coords between index={0}{2}, skip coords between index={3}{5}] table[x=s2sdup, y=s2smiss] {figs/data/fr_scheduler_cal_oneweb_dl.dat};
\addplot[schRing, skip coords between index={0}{2}, skip coords between index={3}{5}] table[x=ds2sdup, y=ds2smiss] {figs/data/fr_scheduler_cal_oneweb_dl.dat};
\draw[black!65, thin] (axis cs:2,2.322) rectangle (axis cs:20,4.984);
\coordinate (zsrcc) at (axis cs:20,3.653);
\nextgroupplot[title={(d) OneWeb UL, $D=583$\,ms}, ymax=18]
\addplot[schOracle] table[x=dup, y=oracle] {figs/data/fr_scheduler_curve_oneweb_ul.dat};
\addplot[schRandom] table[x=dup, y=random] {figs/data/fr_scheduler_curve_oneweb_ul.dat};
\addplot[schReact] table[x=dup, y=reactive] {figs/data/fr_scheduler_curve_oneweb_ul.dat};
\addplot[schS2S] table[x=dup, y=s2s] {figs/data/fr_scheduler_curve_oneweb_ul.dat};
\addplot[schAdapt] table[x=dup, y=s2sadapt] {figs/data/fr_scheduler_curve_oneweb_ul.dat};
\addplot[schGauss] table[x=dup, y=gauss] {figs/data/fr_scheduler_curve_oneweb_ul.dat};
\addplot[schDS2S] table[x=dup, y=ds2s] {figs/data/fr_scheduler_curve_oneweb_ul.dat};
\addplot[schCalS2S] table[x=s2sdup, y=s2smiss] {figs/data/fr_scheduler_cal_oneweb_ul.dat};
\addplot[schCalDS2S] table[x=ds2sdup, y=ds2smiss] {figs/data/fr_scheduler_cal_oneweb_ul.dat};
\addplot[schRing, skip coords between index={0}{2}, skip coords between index={3}{5}] table[x=s2sdup, y=s2smiss] {figs/data/fr_scheduler_cal_oneweb_ul.dat};
\addplot[schRing, skip coords between index={0}{2}, skip coords between index={3}{5}] table[x=ds2sdup, y=ds2smiss] {figs/data/fr_scheduler_cal_oneweb_ul.dat};
\draw[black!65, thin] (axis cs:4,1.729) rectangle (axis cs:16,13.655);
\coordinate (zsrcd) at (axis cs:16,7.692);
\end{groupplot}
\begin{axis}[schInset, name=insa, at={($(group c1r1.north east)+(-0.08cm,-0.08cm)$)}, anchor=north east, xmin=0, xmax=20, ymin=0.688, ymax=7.717]
\addplot[schOracle] table[x=dup, y=oracle] {figs/data/fr_scheduler_curve_starlink_dl.dat};
\addplot[schRandom] table[x=dup, y=random] {figs/data/fr_scheduler_curve_starlink_dl.dat};
\addplot[schReact] table[x=dup, y=reactive] {figs/data/fr_scheduler_curve_starlink_dl.dat};
\addplot[schS2S] table[x=dup, y=s2s] {figs/data/fr_scheduler_curve_starlink_dl.dat};
\addplot[schAdapt] table[x=dup, y=s2sadapt] {figs/data/fr_scheduler_curve_starlink_dl.dat};
\addplot[schGauss] table[x=dup, y=gauss] {figs/data/fr_scheduler_curve_starlink_dl.dat};
\addplot[schDS2S] table[x=dup, y=ds2s] {figs/data/fr_scheduler_curve_starlink_dl.dat};
\addplot[schCalS2S] table[x=s2sdup, y=s2smiss] {figs/data/fr_scheduler_cal_starlink_dl.dat};
\addplot[schCalDS2S] table[x=ds2sdup, y=ds2smiss] {figs/data/fr_scheduler_cal_starlink_dl.dat};
\addplot[schRing, skip coords between index={0}{2}, skip coords between index={3}{5}] table[x=s2sdup, y=s2smiss] {figs/data/fr_scheduler_cal_starlink_dl.dat};
\addplot[schRing, skip coords between index={0}{2}, skip coords between index={3}{5}] table[x=ds2sdup, y=ds2smiss] {figs/data/fr_scheduler_cal_starlink_dl.dat};
\end{axis}
\draw[-latex, black!65, thin] (zsrca) -- (insa.south west);
\begin{axis}[schInset, name=insb, at={($(group c2r1.north east)+(-0.08cm,-0.08cm)$)}, anchor=north east, xmin=2, xmax=20, ymin=1.663, ymax=4.212]
\addplot[schOracle] table[x=dup, y=oracle] {figs/data/fr_scheduler_curve_starlink_ul.dat};
\addplot[schRandom] table[x=dup, y=random] {figs/data/fr_scheduler_curve_starlink_ul.dat};
\addplot[schReact] table[x=dup, y=reactive] {figs/data/fr_scheduler_curve_starlink_ul.dat};
\addplot[schS2S] table[x=dup, y=s2s] {figs/data/fr_scheduler_curve_starlink_ul.dat};
\addplot[schAdapt] table[x=dup, y=s2sadapt] {figs/data/fr_scheduler_curve_starlink_ul.dat};
\addplot[schGauss] table[x=dup, y=gauss] {figs/data/fr_scheduler_curve_starlink_ul.dat};
\addplot[schDS2S] table[x=dup, y=ds2s] {figs/data/fr_scheduler_curve_starlink_ul.dat};
\addplot[schCalS2S] table[x=s2sdup, y=s2smiss] {figs/data/fr_scheduler_cal_starlink_ul.dat};
\addplot[schCalDS2S] table[x=ds2sdup, y=ds2smiss] {figs/data/fr_scheduler_cal_starlink_ul.dat};
\addplot[schRing, skip coords between index={0}{2}, skip coords between index={3}{5}] table[x=s2sdup, y=s2smiss] {figs/data/fr_scheduler_cal_starlink_ul.dat};
\addplot[schRing, skip coords between index={0}{2}, skip coords between index={3}{5}] table[x=ds2sdup, y=ds2smiss] {figs/data/fr_scheduler_cal_starlink_ul.dat};
\end{axis}
\draw[-latex, black!65, thin] (zsrcb) -- (insb.south west);
\begin{axis}[schInset, name=insc, at={($(group c3r1.north east)+(-0.08cm,-0.08cm)$)}, anchor=north east, xmin=2, xmax=20, ymin=2.322, ymax=4.984]
\addplot[schOracle] table[x=dup, y=oracle] {figs/data/fr_scheduler_curve_oneweb_dl.dat};
\addplot[schRandom] table[x=dup, y=random] {figs/data/fr_scheduler_curve_oneweb_dl.dat};
\addplot[schReact] table[x=dup, y=reactive] {figs/data/fr_scheduler_curve_oneweb_dl.dat};
\addplot[schS2S] table[x=dup, y=s2s] {figs/data/fr_scheduler_curve_oneweb_dl.dat};
\addplot[schAdapt] table[x=dup, y=s2sadapt] {figs/data/fr_scheduler_curve_oneweb_dl.dat};
\addplot[schGauss] table[x=dup, y=gauss] {figs/data/fr_scheduler_curve_oneweb_dl.dat};
\addplot[schDS2S] table[x=dup, y=ds2s] {figs/data/fr_scheduler_curve_oneweb_dl.dat};
\addplot[schCalS2S] table[x=s2sdup, y=s2smiss] {figs/data/fr_scheduler_cal_oneweb_dl.dat};
\addplot[schCalDS2S] table[x=ds2sdup, y=ds2smiss] {figs/data/fr_scheduler_cal_oneweb_dl.dat};
\addplot[schRing, skip coords between index={0}{2}, skip coords between index={3}{5}] table[x=s2sdup, y=s2smiss] {figs/data/fr_scheduler_cal_oneweb_dl.dat};
\addplot[schRing, skip coords between index={0}{2}, skip coords between index={3}{5}] table[x=ds2sdup, y=ds2smiss] {figs/data/fr_scheduler_cal_oneweb_dl.dat};
\end{axis}
\draw[-latex, black!65, thin] (zsrcc) -- (insc.south west);
\begin{axis}[schInset, name=insd, at={($(group c4r1.north east)+(-0.08cm,-0.08cm)$)}, anchor=north east, xmin=4, xmax=16, ymin=1.729, ymax=13.655]
\addplot[schOracle] table[x=dup, y=oracle] {figs/data/fr_scheduler_curve_oneweb_ul.dat};
\addplot[schRandom] table[x=dup, y=random] {figs/data/fr_scheduler_curve_oneweb_ul.dat};
\addplot[schReact] table[x=dup, y=reactive] {figs/data/fr_scheduler_curve_oneweb_ul.dat};
\addplot[schS2S] table[x=dup, y=s2s] {figs/data/fr_scheduler_curve_oneweb_ul.dat};
\addplot[schAdapt] table[x=dup, y=s2sadapt] {figs/data/fr_scheduler_curve_oneweb_ul.dat};
\addplot[schGauss] table[x=dup, y=gauss] {figs/data/fr_scheduler_curve_oneweb_ul.dat};
\addplot[schDS2S] table[x=dup, y=ds2s] {figs/data/fr_scheduler_curve_oneweb_ul.dat};
\addplot[schCalS2S] table[x=s2sdup, y=s2smiss] {figs/data/fr_scheduler_cal_oneweb_ul.dat};
\addplot[schCalDS2S] table[x=ds2sdup, y=ds2smiss] {figs/data/fr_scheduler_cal_oneweb_ul.dat};
\addplot[schRing, skip coords between index={0}{2}, skip coords between index={3}{5}] table[x=s2sdup, y=s2smiss] {figs/data/fr_scheduler_cal_oneweb_ul.dat};
\addplot[schRing, skip coords between index={0}{2}, skip coords between index={3}{5}] table[x=ds2sdup, y=ds2smiss] {figs/data/fr_scheduler_cal_oneweb_ul.dat};
\end{axis}
\draw[-latex, black!65, thin] (zsrcd) -- (insd.south west);
\path ($(group c1r1.north west)!0.5!(group c4r1.north east)$) ++(0,0.62cm) node[anchor=south] {\pgfplotslegendfromname{legsched}};
\end{tikzpicture}
}
  \caption{Deadline-aware duplication on the test sessions: share of late, non-duplicated packets against the share of duplicated packets, pooled over 0.2--3\,s ($D$: 90th percentile of the calibration-session RTT). Curves sweep the knob of each policy. Markers: one-sided bounds calibrated at levels of 0.80 to 0.99, circled at 0.95. Insets: low duplication. S2S-E adapt.\ scales its margin by the input volatility.}
  \label{fig:sched}
\end{figure*}

Each policy has one knob, and sweeping it traces a curve in Fig.~\ref{fig:sched}. A deployed scheduler must set the knob in advance. For S2S-E we use the one-sided 95\% conformal upper bound calibrated on the calibration sessions. For DS2S-E we set the knob on the calibration sessions so that its miss rate there equals that of the S2S-E bound. On Starlink DL, the S2S-E bound duplicates \SchedSSdupSLDL\% of the packets and misses \SchedSSmissSLDL\%. DS2S-E duplicates \SchedDSdupSLDL\% and misses \SchedDSmissSLDL\%, which is \SchedSaveSLDL\% less duplication (95\% CI \SchedSaveSLDLlo{} to \SchedSaveSLDLhi\%) with no more misses. On Starlink UL and OneWeb UL the two schedulers tie (duplication change \SchedSaveSLUL\% and \SchedSaveOWUL\%). On OneWeb DL, DS2S-E needs \SchedDSdupCurveOWDL\% duplication to reach the miss rate of the S2S-E bound, which duplicates \SchedSSdupOWDL\%, and with its knob set on the calibration sessions it misses \SchedDSmissOWDL\% of the packets against \SchedSSmissOWDL\%. Every forecaster with the anchored decoder loses there, including the point forecaster S2S-AE-E. The loss therefore comes from the decoder family and not from the distributional output.

On Starlink DL the gain has two parts. With the same inputs and decoder, the point forecaster S2S-AE-E needs \SchedFFdupSLDL\% duplication at the miss rate of the S2S-E bound, and DS2S-E needs \SchedDSdupCurveSLDL\% (S2S-E: \SchedSSdupCurveSLDL\%). About half of the gain thus comes from the new inputs and decoder and half from the predictive distribution. With both knobs set on the calibration sessions, DS2S-E needs \SchedSaveFFSLDL\% less duplication than S2S-AE-E. \GaussSchedSentence{} Fig.~\ref{fig:risk} shows what a point forecast misses. It groups the test windows by the S2S-E mean forecast and, within each group, by the spread that DS2S-E predicts. At the same mean forecast of \RiskMeanLo{} to \RiskMeanHi\,ms on Starlink DL, \RiskLowPct\% of the packets are late when DS2S-E predicts a narrow spread and \RiskHighPct\% when it predicts a wide one. A scheduler that sees only the mean must treat these windows alike.

\begin{figure*}[t]
  \centering
  % fr_scheduler_risk.tex - generated by modeling/experiments/fr_scheduler.py (edit the script, not this file).
% Test deadline-violation rate (D = validation p90) per S2S-E mean bin (x) and DS2S-E spread quintile
% inside that bin (y). Gray: fewer than 200 decisions.
\begin{tikzpicture}
\begin{groupplot}[group style={group size=4 by 1, horizontal sep=0.45cm, y descriptions at=edge left},
  frbase, scale only axis, width=3.3cm, height=2.0cm, grid=none, enlargelimits=false,
  axis background/.style={fill=gray!25}, tick style={draw=none},
  xmin=0.5, xmax=10.5, ymin=0.5, ymax=5.5, ytick={1,3,5}, yticklabels={low,mid,high},
  xlabel={S2S-E mean forecast (ms)}, ylabel={DS2S-E spread},
  colormap/viridis, point meta min=-1, point meta max=2,
  title style={font=\footnotesize, yshift=-0.6ex}]
\nextgroupplot[title={(a) Starlink DL}, xtick={1,5,10}, xticklabels={38,92,189}]
\addplot[matrix plot*, mesh/cols=10, point meta=explicit] table[x=mean_bin, y=spread_q, meta=lrate] {figs/data/fr_scheduler_risk_starlink_dl.dat};
\nextgroupplot[title={(b) Starlink UL}, xtick={1,5,10}, xticklabels={32,44,50}]
\addplot[matrix plot*, mesh/cols=10, point meta=explicit] table[x=mean_bin, y=spread_q, meta=lrate] {figs/data/fr_scheduler_risk_starlink_ul.dat};
\nextgroupplot[title={(c) OneWeb DL}, xtick={1,5,10}, xticklabels={58,63,108}]
\addplot[matrix plot*, mesh/cols=10, point meta=explicit] table[x=mean_bin, y=spread_q, meta=lrate] {figs/data/fr_scheduler_risk_oneweb_dl.dat};
\nextgroupplot[title={(d) OneWeb UL}, xtick={1,5,10}, xticklabels={89,102,779}, colorbar, colorbar style={width=0.18cm, tick label style={font=\scriptsize}, ytick={-1,0,1,2}, yticklabels={0.1,1,10,100}, ylabel={late packets (\%)}, ylabel style={font=\scriptsize}}]
\addplot[matrix plot*, mesh/cols=10, point meta=explicit] table[x=mean_bin, y=spread_q, meta=lrate] {figs/data/fr_scheduler_risk_oneweb_ul.dat};
\end{groupplot}
\end{tikzpicture}

  \caption{Same mean, different risk. Share of late packets ($D$ as in Fig.~\ref{fig:sched}, log color scale) when the test windows are grouped by the S2S-E mean forecast (deciles, x) and, within each decile, by the predicted DS2S-E spread $\hat{q}_{0.975}-\hat{q}_{0.5}$ (quintiles, y).}
  \label{fig:risk}
\end{figure*}

\subsection{Sensitivity, Other Targets, and Cost}
\label{sec:sens}
The advantage of DS2S over S2S holds for 3\,s and 12\,s input windows and for a 6\,s horizon on the four satellite links (two seeds each). The interval score improves by \SweepISLoAbs\% to \SweepISHiAbs\% and the MAE improves on every link in every setting, while the RMSE changes by \SweepRMSELo\% to $+$\SweepRMSEHi\%. \WeightSentence{} For jitter and loss, which are one-second iperf reports with long runs of zeros, the point accuracy of DS2S and S2S is similar, and the interval score of DS2S is lower for both targets on all five links where they vary, by up to a factor of \JitISRatio{} (Starlink DL jitter).

What does DS2S use? Its forecast is most sensitive to the newest second of input, which receives \GradLastLo\% to \GradLastHi\% of the input gradient of the learned change (17\% if uniform). Its attention weights look elsewhere. Replacing them by uniform weights raises the RTT MAE by at most \AttnAblMax\% on Starlink and 5G and matters only on OneWeb (OneWeb DL: $+$\AttnAblOWDL\%). The gain therefore does not come from where the model attends. On one CPU thread, a forecast of all 15 steps and three targets takes \TimeDSCPU\,ms, far below the 200\,ms between measurements.
